\section{DESA}
\label{sec:method}

\subsection{Coupled Construction Overview}

DESA jointly designs dense and sparse query construction for the fused result,
aiming to improve complete-list effectiveness while reducing both certification
depths relative to the unexpanded query. Given an original query $q_o$, DESA
first generates one set of complementary reference passages
$R=\{r_1,\ldots,r_n\}$ shared by both channels. This common evidence basis
prevents channel differences from originating in independently generated
rewrites, placing specialization entirely in integration. Dense retrieval may
incorporate new semantic directions from the references, whereas sparse
retrieval may only reorder documents supported by the original query. The
resulting dense and sparse rankings are fused into the final result.
Figure~\ref{fig:desa-overview} summarizes this coupled construction.

\begin{figure}[H]
  \centering
  \includegraphics[width=\columnwidth]{figures/fig_desa_overview.pdf}
  \caption{Coupled construction in DESA. A shared reference set establishes a
  common evidence basis. Orthogonal residual expansion adds new dense semantic
  directions, whereas score-product anchoring restricts sparse reordering to
  the original query's support.}
  \label{fig:desa-overview}
\end{figure}

\subsection{Shared Reference Generation}

DESA generates one reference set before constructing either channel ranking,
rather than producing a separate rewrite for each retriever. The generated
evidence is therefore held common across channels, so their differences arise
from controlled integration rather than different generated content.

The generator is prompted to produce $n$ short passages that could plausibly
occur in relevant documents. Each passage introduces a different alias,
technical term, or paraphrase while retaining the entities, relations,
numbers, polarity, temporal conditions, scope, and other constraints of
$q_o$. The prompt prohibits broader information needs and invented answers.
The passages therefore vary in expression rather than intent. The same set
$R$ is supplied unchanged to both channel operators. The complete prompt and
decoding settings are provided in Appendix~A.

\subsection{Coupled Dense--Sparse Construction}

The two operators act on the same original query and generated references but
assign different roles to the shared evidence. Dense retrieval permits the
references to add semantic information, subject to bounded deviation from the
original query. Sparse retrieval permits them to alter ranking order, but not
to admit documents unsupported by the original query.

\subsubsection{Orthogonal Residual Expansion}

The dense branch should incorporate semantic information absent from the
original query without allowing the generated references to replace its
original direction. Let $f_D$ be the dense encoder. We encode the original
query and each query--reference pair as unit vectors:

\[
\begin{aligned}
\mathbf v_o&=\operatorname{norm}(f_D(q_o)),\\
\mathbf v_i&=\operatorname{norm}
\!\left(f_D(\operatorname{concat}(q_o,r_i))\right),\\
\bar{\mathbf v}_r&=\frac{1}{n}\sum_{i=1}^{n}\mathbf v_i .
\end{aligned}
\]

We remove the component of $\bar{\mathbf v}_r$ parallel to the original query
and add the remaining residual to $\mathbf v_o$:

\[
\mathbf z=
\bar{\mathbf v}_r-
(\mathbf v_o^\top\bar{\mathbf v}_r)\mathbf v_o,
\qquad
\mathbf e_D=\operatorname{norm}(\mathbf v_o+\mathbf z).
\]

Documents are ranked by their similarity to $\mathbf e_D$, producing the
dense ranking $\pi_D$. Because $\mathbf z\perp\mathbf v_o$ and
$\|\mathbf z\|\leq1$,

\[
\angle(\mathbf e_D,\mathbf v_o)
=\arctan\|\mathbf z\|
\leq\frac{\pi}{4}.
\]

The references can therefore add semantic directions absent from the original
representation while keeping the expanded query within $45^\circ$ of it. If
they contribute no orthogonal component, then $\mathbf e_D=\mathbf v_o$.
Orthogonal residual expansion thus grants the shared evidence additive
semantic authority through a bounded, parameter-free update.

\subsubsection{Score-Product Anchoring}

The sparse branch requires a different constraint. Directly adding generated
terms may admit expansion-only documents, changing the sparse support and its
relation to the dense ranking. DESA instead allows generated lexical evidence
to reorder documents only within the original sparse support.

The expanded query concatenates the original query and all reference passages:

\[
q_r=\operatorname{concat}(q_o,r_1,\ldots,r_n).
\]

Let $s_S(q,d)$ be a non-negative additive sparse scorer. Each document $d$
receives an original score, an expanded score, and their product:

\[
\begin{aligned}
s_o(d)&=s_S(q_o,d), &
s_r(d)&=s_S(q_r,d),\\
s_A(d)&=s_o(d)s_r(d).&&
\end{aligned}
\]

Because $q_r$ contains $q_o$, every document supported by the original query
is also supported by the expanded query. With
$\operatorname{supp}(s)=\{d:s(d)>0\}$,

\[
\operatorname{supp}(s_A)
=\operatorname{supp}(s_o)\cap\operatorname{supp}(s_r)
=\operatorname{supp}(s_o).
\]

Sorting the positive anchored scores produces the sparse ranking $\pi_S$.
Generated terms may promote or demote documents within the original support,
but documents matched only by the expansion receive zero anchored score and
are excluded. If $s_r$ is merely a positive multiple of $s_o$, the original
sparse order is unchanged; positive global score rescaling also leaves the
ranking invariant. Score-product anchoring therefore grants the shared
evidence reordering authority but no independent sparse admission authority.
